%% ma: L10-B13-0103
%% lop: 10
%% chuong: 5
%% bai: 13
%% dang: 10.13.13
%% loai: DS
%% muc: [NB, TH, TH, TH]
%% dap_an: "SSĐĐ"
%% nguon: "chuyen-de-cac-so-dac-trung-cua-mau-so-lieu-khong-ghep-nhom-toan-10-chuong-trinh-moi.pdf (GV. Phan Nhật Linh) (P03), Bài 02 Dạng 2 Phần II Câu 7"
%% kien_thuc: "Lớp 10 Bài 13 (số trung bình, trung vị, tứ phân vị, mốt)"
%% trang_thai: chinh-thuc
%% ghi_chu: "Tự duyệt theo yêu cầu giáo viên 22/10/2026."
\begin{ex}
Nhiệt độ trung bình ($^\circ$C) hằng tháng trong một năm được ghi lại trong bảng sau:
\begin{center}
\begin{tabular}{|c|c|c|c|c|c|c|c|c|c|c|c|c|}\hline
Tháng & $1$ & $2$ & $3$ & $4$ & $5$ & $6$ & $7$ & $8$ & $9$ & $10$ & $11$ & $12$\\\hline
Nhiệt độ & $16$ & $20$ & $25$ & $28$ & $30$ & $30$ & $28$ & $25$ & $25$ & $20$ & $18$ & $16$\\\hline
\end{tabular}
\end{center}
Xét tính đúng sai của các khẳng định sau:
\choiceTF{Mốt của mẫu số liệu là $30$}{Nhiệt độ trung bình của các tháng trong năm là $25\,^\circ$C}{\True Trung vị của mẫu số liệu là $25$}{\True Tứ phân vị thứ ba của mẫu số liệu là $28$}
\loigiai{Sắp xếp: $16$; $16$; $18$; $20$; $20$; $25$; $25$; $25$; $28$; $28$; $30$; $30$.

a) Sai: $25$ xuất hiện $3$ lần, nhiều nhất, nên $M_o=25$.

b) Sai: $\overline x=\dfrac{281}{12}\approx23{,}42$.

c) Đúng: $M_e=\dfrac{25+25}2=25$.

d) Đúng: $Q_3=\dfrac{x_9+x_{10}}2=\dfrac{28+28}2=28$.}
\end{ex}
